% \begin{document}

In proofs, it will often be convenient to have structural rules that allow us to rearrange
assertions in various ways. The \rn{Exists} rule lets us reason about a command under an
existential quantifier: if $c$ takes $H$ to $J$ for every value of the logical variable $i$, then
it takes $\exists i.\,H$ to $\exists i.\,J$.
\[
  \infr{\vdash\triple{H}{c}{J}}{\vdash\triple{\exists i.\,H}{c}{\exists i.\,J}}{Exists}
\]
In the examples, we will use this rule constantly, because unfolding a recursive predicate will
always introduces existentials.

\section{Example: Disposing a List}

The following program $\mathrm{dispose}$ frees every node of the list at $p$:
\[
  \begin{array}{l}
    \kw{while}\ (p\neq 0)\ \kw{do}\ (\\
    \quad q:=\ast(p+1);\\
    \quad\kw{free}(p);\\
    \quad\kw{free}(p+1);\\
    \quad p:=q\\
    )
  \end{array}
\]
Its specification is $\{\lst(p,L)\}\ \mathrm{dispose}\ \{\emp\}$, and a suitable loop invariant is
$\exists L'.\,\lst(p,L')$: the heap always holds exactly the rest of the list. The precondition
implies the invariant (take $L'=L$). For the loop body, the \rn{While} rule requires us to show
$\{\pure{p\neq 0}\sep\exists L'.\,\lst(p,L')\}\ \mathrm{body}\ \{\exists L'.\,\lst(p,L')\}$:
\begin{align*}
  &\{\,(\exists L'.\,\lst(p,L'))\sep\pure{p\neq 0}\,\}\ \Rightarrow\ \text{(unfold)}\\
  &\{\,\exists v,y,L''.\ p\pto v\sep(p+1)\pto y\sep\lst(y,L'')\,\}\\
  &q:=\ast(p+1);\qquad\text{(\rn{Load}, \rn{Frame})}\\
  &\{\,\exists v,y,L''.\ \pure{q=y}\sep p\pto v\sep(p+1)\pto y\sep\lst(y,L'')\,\}\\
  &\kw{free}(p);\qquad\text{(\rn{Free}, \rn{Frame})}\\
  &\{\,\exists y,L''.\ \pure{q=y}\sep(p+1)\pto y\sep\lst(y,L'')\,\}\\
  &\kw{free}(p+1);\qquad\text{(\rn{Free}, \rn{Frame})}\\
  &\{\,\exists y,L''.\ \pure{q=y}\sep\lst(y,L'')\,\}\ \Rightarrow\\
  &\{\,\exists L'.\,\lst(q,L')\,\}\\
  &p:=q\qquad\text{(assignment)}\\
  &\{\,\exists L'.\,\lst(p,L')\,\}
\end{align*}
The first implication unfolds the list: since $p\neq 0$, the list cannot be empty (as
$\lst(p,[\,])$ would force $p=0$), so $L'=v\ld L''$ for some $v$ and $L''$. Each command is then
verified under the existential using \rn{Exists}, and with \rn{Frame} to carry along the cells it
does not touch. For example, the load uses the \rn{Load} axiom on the single cell $(p+1)\pto y$ and
frames $p\pto v\sep\lst(y,L'')$; its side condition holds because $q$ occurs in neither $p+1$ nor
$y$. After $\kw{free}(p)$, no cell mentions $v$ any more, so we drop $\exists v$ by
\rn{Consequence}. The last implication repackages what remains as the invariant with $q$ in place
of $p$, and the assignment axiom then gives the invariant itself.

When the loop exits, we have $(\exists L'.\,\lst(p,L'))\sep\pure{p=0}$. Since $0\pto v$ is
unsatisfiable, $\lst(0,L')$ forces $L'=[\,]$, and so the heap is empty, as required. Because
triples are partial correctness statements, we did not need to prove that the loop terminates
(although it does, since the length of $L'$ decreases on every iteration).

\section{List Segments}

For loops that traverse a list without destroying it, we need to describe the part of the list that
has already been visited. A \emph{list segment} from $p$ to $q$ holding $L$ is a list that ends at
$q$ instead of at $0$:
\begin{align*}
  \lseg(p,q,[\,])&\ \triangleq\ \pure{p=q}\\
  \lseg(p,q,v\ld L)&\ \triangleq\ \exists r.\ p\pto v\sep(p+1)\pto r\sep\lseg(r,q,L)
\end{align*}

\begin{center}
\begin{tikzpicture}[>=Stealth,
    cell/.style={draw,minimum width=8mm,minimum height=7mm,inner sep=0pt}]
  \node (p) at (-1,0) {$p$};
  \draw[->] (p) -- (-0.45,0);
  \foreach \i/\v in {0/v_1,2/v_2,4/v_3} {
    \node[cell] (v\i) at (\i,0) {$\v$};
    \node[cell] (n\i) at (\i+0.8,0) {$\bullet$};
  }
  \draw[->] (n0.center) -- (v2.west);
  \draw[->] (n2.center) -- (v4.west);
  \draw[->] (n4.center) -- ++(1.0,0) node[right] {$q$};
  \node at (3,-1) {$\lseg(p,q,[v_1,v_2,v_3])$};
\end{tikzpicture}
\end{center}

The following facts are each proved by induction on $L$:
\begin{align*}
  \lst(p,L)&\equiv\lseg(p,0,L)\\
  \lseg(p,x,L)\sep x\pto v\sep(x+1)\pto y&\Rightarrow\lseg(p,y,L\app[v]) && \text{(snoc)}\\
  \lseg(p,y,L)\sep\lst(y,L')&\Rightarrow\lst(p,L\app L') && \text{(concat)}
\end{align*}
Snoc extends a segment by one node at its end, and concat attaches a list to the end of a segment.

\section{Example: List Append}

The following program $\mathrm{append}$, appends the list at $p_2$ to the list at $p_1$, leaving
the result at $r$.
\[
  \begin{array}{l}
    x:=p_1;\\
    n:=\ast(x+1);\\
    \kw{while}\ (n\neq 0)\ \kw{do}\ (\\
    \quad x:=n;\\
    \quad n:=\ast(x+1)\\
    );\\
    \ast(x+1):=p_2;\\
    r:=p_1
  \end{array}
\]
In the program, the variable $x$ always points at a node of the first list, and $n$ holds that
node's next pointer; the loop walks $x$ to the last node, and the final store overwrites its null
pointer with $p_2$.

For example, suppose the heap initially holds the lists $[1,2]$ at address $6$ and $[3,4]$ at
address $99$:
\begin{center}
\begin{tabular}{l|cccccccc}
  address & 6 & 7 & 22 & 23 & 46 & 47 & 99 & 100\\\hline
  before  & 1 & 46 & 4 & 0 & 2 & 0 & 3 & 22\\
  after   & 1 & 46 & 4 & 0 & 2 & 99 & 3 & 22
\end{tabular}
\end{center}
with $p_1=6$ and $p_2=99$. Afterwards, $x=46$ (the last node of the first list), $n=0$, $r=6$, and
the heap holds the list $[1,2,3,4]$ at $r$. Only one cell changes---the last node is modified from
$0$ to $p_2$.

The specification for this program
\[
  \{\lst(p_1,L_1)\sep\lst(p_2,L_2)\sep\pure{p_1\neq 0}\}\ \mathrm{append}\ \{\lst(r,L_1\app L_2)\}.
\]
Two aspects of this specification deserve comment. First, the pure fact $p_1\neq 0$ is essential:
if the first list is empty, then $x=0$, so the program accesses address $1$, which need not be
allocated, and it may produce a memory error. Second, $\lst(p_2,L_2)$ does not appear in the
postcondition because its cells have become part of $\lst(r,L_1\app L_2)$. The precondition's
$\sep$ guarantees that the two lists are disjoint, so appending cannot create a cycle.

To keep the proof readable, we abbreviate a node as
\[
  \nd{p}{v}{q}\ \triangleq\ p\pto v\sep(p+1)\pto q.
\]
The proof has three parts: the code before the loop, which establishes the invariant; the loop
body, which preserves it; and the code after the loop, which establishes the postcondition.

\paragraph{Before the loop.} We frame $\lst(p_2,L_2)$, which this code does not touch. Since
$p_1\neq 0$, the first list is not empty, so we can unfold its first node:
\begin{align*}
  &\{\,\lst(p_1,L_1)\sep\pure{p_1\neq 0}\,\}\ \Rightarrow\\
  &\{\,\exists v,y,L_b.\ \pure{L_1=v\ld L_b}\sep\nd{p_1}{v}{y}\sep\lst(y,L_b)\,\}\\
  &x:=p_1;\\
  &\{\,\exists v,y,L_b.\ \pure{L_1=v\ld L_b}\sep\pure{x=p_1}\sep\nd{x}{v}{y}\sep\lst(y,L_b)\,\}\\
  &n:=\ast(x+1);\\
  &\{\,\exists v,y,L_b.\ \pure{L_1=v\ld L_b}\sep\pure{x=p_1}\sep\pure{n=y}\sep\nd{x}{v}{y}\sep\lst(y,L_b)\,\}\ \Rightarrow\\
  &\{\,\exists v,L_b.\ \pure{L_1=[\,]\app v\ld L_b}\sep\lseg(p_1,x,[\,])\sep\nd{x}{v}{n}\sep\lst(n,L_b)\,\}
\end{align*}
Note that the pure fact $\pure{L_1=v\ld L_b}$ must be carried through every assertion: it is the
only connection between the logical variables $v$ and $L_b$ and the list $L_1$ in the
specification. In the last step we use $\lseg(p_1,x,[\,])\equiv\pure{p_1=x}$ and replace $y$ by
$n$. Generalizing $[\,]$ to an arbitrary $L_a$, and adding back the frame, suggests the loop
invariant
\[
  I\ \triangleq\ \exists L_a,v,L_b.\ \pure{L_1=L_a\app v\ld L_b}\sep\lseg(p_1,x,L_a)\sep\nd{x}{v}{n}\sep\lst(n,L_b)\sep\lst(p_2,L_2).
\]
The visited part of the first list is the segment from $p_1$ to $x$; $x$ is the current node,
holding $v$ and pointing to $n$; and the rest of the first list is at $n$. The assertion we derived
is exactly $I$ with $L_a=[\,]$, so the code before the loop establishes the invariant.

\paragraph{The loop body.} We must show
$\{I\sep\pure{n\neq 0}\}\ x:=n;\ n:=\ast(x+1)\ \{I\}$. We again frame $\lst(p_2,L_2)$. Since
$n\neq 0$, the rest of the list is not empty, so we unfold $L_b=w\ld L_c$:
\begin{align*}
  &\{\,I\sep\pure{n\neq 0}\,\}\ \Rightarrow\ \text{(unfold)}\\
  &\{\,\exists L_a,v,w,m,L_c.\ \pure{L_1=L_a\app v\ld w\ld L_c}\\
  &\qquad\sep\lseg(p_1,x,L_a)\sep\nd{x}{v}{n}\sep\nd{n}{w}{m}\sep\lst(m,L_c)\,\}\ \Rightarrow\ \text{(snoc)}\\
  &\{\,\exists L_a,v,w,m,L_c.\ \pure{L_1=L_a\app v\ld w\ld L_c}\\
  &\qquad\sep\lseg(p_1,n,L_a\app[v])\sep\nd{n}{w}{m}\sep\lst(m,L_c)\,\}\\
  &x:=n;\qquad\text{(assignment)}\\
  &\{\,\exists L_a,v,w,m,L_c.\ \pure{L_1=L_a\app v\ld w\ld L_c}\\
  &\qquad\sep\lseg(p_1,x,L_a\app[v])\sep\nd{x}{w}{m}\sep\lst(m,L_c)\,\}\\
  &n:=\ast(x+1)\qquad\text{(\rn{Load}, \rn{Frame})}\\
  &\{\,\exists L_a,v,w,m,L_c.\ \pure{L_1=L_a\app v\ld w\ld L_c}\sep\pure{n=m}\\
  &\qquad\sep\lseg(p_1,x,L_a\app[v])\sep\nd{x}{w}{m}\sep\lst(m,L_c)\,\}\ \Rightarrow\\
  &\{\,I\,\}
\end{align*}
The snoc step folds the current node into the visited segment. The final implication re-establishes
$I$ with $L_a\app[v]$ for $L_a$, $w$ for $v$, and $L_c$ for $L_b$. This fact holds because
$L_a\app v\ld w\ld L_c=(L_a\app[v])\app w\ld L_c$, and $\nd{x}{w}{m}\sep\lst(m,L_c)$ becomes
$\nd{x}{w}{n}\sep\lst(n,L_c)$ because $n=m$.

\paragraph{After the loop.} When the loop exits, we have $I\sep\pure{n=0}$. Since $n=0$, the rest
of the list must be empty, so $L_b=[\,]$ and hence $L_1=L_a\app[v]$:
\begin{align*}
  &\{\,I\sep\pure{n=0}\,\}\ \Rightarrow\\
  &\{\,\exists L_a,v.\ \pure{L_1=L_a\app[v]}\sep\lseg(p_1,x,L_a)\sep\nd{x}{v}{0}\sep\lst(p_2,L_2)\,\}\\
  &\ast(x+1):=p_2;\qquad\text{(\rn{Store}, \rn{Frame})}\\
  &\{\,\exists L_a,v.\ \pure{L_1=L_a\app[v]}\sep\lseg(p_1,x,L_a)\sep\nd{x}{v}{p_2}\sep\lst(p_2,L_2)\,\}\ \Rightarrow\ \text{(snoc)}\\
  &\{\,\lseg(p_1,p_2,L_1)\sep\lst(p_2,L_2)\,\}\\
  &r:=p_1\qquad\text{(assignment)}\\
  &\{\,\lseg(r,p_2,L_1)\sep\lst(p_2,L_2)\,\}\ \Rightarrow\ \text{(concat)}\\
  &\{\,\lst(r,L_1\app L_2)\,\}
\end{align*}
This establishes the postcondition and completes the proof.

\section{The Magic Wand}

In standard Hoare logic, the assignment axiom $\{P[e/x]\}\ x:=e\ \{P\}$ works backwards: it
computes the weakest precondition of an assignment by substitution. What is the weakest
precondition of a heap store such as $\ast x:=5$ with respect to a postcondition $Q$? Substitution
doesn't work, because $Q$ may refer to the cell at $x$ under a different name---e.g., an assertion
such as $y\pto 3$ may refers to the same cell when $x$ and $y$ are aliases. What we would like to
say is: ``I hold the cell at $x$, together with a promise: add $x\pto 5$, and you get $Q$.'' Our
existing assertion language has no way to express such a promise, so we add a new connective,
called \emph{magic wand}.

\paragraph{Definition.} The \emph{magic wand}, or \emph{separating implication}, is defined by
\[
  (\sigma,h)\models_I H\wand J\quad\text{iff for all } h'\perp h \text{ with } (\sigma,h')\models_I H,
  \text{ we have } (\sigma,h\uplus h')\models_I J.
\]
Intuitively, $H\wand J$ describes a heap that becomes a $J$-heap when you add an $H$-heap; it is
sometimes described as ``$J$ with an $H$-shaped hole.'' The wand is to $\sep$ as implication is to
conjunction. In particular, it satisfies modus ponens and a currying law:
\[
  H\sep(H\wand J)\Rightarrow J
  \qquad\qquad
  (K\sep H\Rightarrow J)\ \text{iff}\ (K\Rightarrow H\wand J)
\]
For example, $\emp\wand H\equiv H$, and any heap satisfying $y\pto 4$ also satisfies
\[
  x\pto 3\wand(x\pto 3\sep y\pto 4).
\]
(If $x$ and $y$ are aliases, no disjoint heap satisfies $x\pto 3$, so the wand holds vacuously.)

\paragraph{Weakest preconditions for heap commands.} With the wand, we can state the weakest
preconditions of the heap commands, where ``weakest'' is with respect to our fault-avoiding
triples:
\begin{center}
\begin{tabular}{ll}
  \textbf{Command} & \textbf{Weakest precondition for $Q$}\\\hline
  $\ast e_1:=e_2$ & $(e_1\pto-)\sep((e_1\pto e_2)\wand Q)$\\
  $\kw{free}(e)$ & $(e\pto-)\sep Q$\\
  $x:=\kw{new}(e)$ & $\forall i.\ (i\pto e)\wand Q[i/x]$\quad($i$ fresh)
\end{tabular}
\end{center}
For a store, we must own the cell at $e_1$ (so that the store does not fault), and the rest of the
heap must turn into a $Q$-heap once a cell $e_1\pto e_2$ is added. For allocation, the rest of the
heap must turn into a $Q$-heap once any new cell $i\pto e$ is added, with $x$ holding its address
$i$, because \rn{New} may choose any unused address. For example, with
$Q=x\pto 5\sep y\pto 7$,
\[
  \wpre(\ast x:=5,\,Q)=x\pto-\sep(x\pto 5\wand(x\pto 5\sep y\pto 7))\ \Leftarrow\ x\pto-\sep y\pto 7,
\]
by currying, since $y\pto 7\sep x\pto 5\Rightarrow x\pto 5\sep y\pto 7$. This is exactly what we
would expect.

\paragraph{Backward rules from forward ones.} We do not need new proof rules to use these weakest
preconditions. For example, the backward rule for stores follows from the forward \rn{Store}
axiom. Let $K=(e_1\pto e_2)\wand Q$. Since $\modv(\ast e_1:=e_2)=\emptyset$, the \rn{Store} axiom
and \rn{Frame} give
\[
  \{e_1\pto-\sep K\}\ \ast e_1:=e_2\ \{e_1\pto e_2\sep K\}.
\]
By modus ponens, $e_1\pto e_2\sep((e_1\pto e_2)\wand Q)\Rightarrow Q$, so by \rn{Consequence},
\[
  \{(e_1\pto-)\sep((e_1\pto e_2)\wand Q)\}\ \ast e_1:=e_2\ \{Q\}.
\]
In other words, the forward small axioms together with \rn{Frame} already suffice; the wand simply
lets us state weakest preconditions.

\paragraph{Where the magic wand shows up.} The wand plays a central role in modern program logics.
In Iris, for example, the rules for weakest preconditions are stated using $\wand$, much like the
backward rules above. Viper supports magic wands in specifications. The wand is also useful for
describing partial data structures: for instance, $\lst(x,L_2)\wand\lst(p,L_1\app L_2)$ describes
a list with a hole, which becomes a complete list once the list at $x$ is supplied. On the other
hand, the wand is notoriously hard to automate, which is one reason why most tools reason in the
forward direction, often using symbolic execution.

\section*{Further Reading}

Reynolds's lecture notes cover list segments, the structural rules (including the form of the
\rn{Exists} rule used here), and backward reasoning with the magic wand. Chargu\'eraud's textbook
\emph{Separation Logic Foundations} (Software Foundations, Volume 6) takes as primitive a variant of
the \rn{Exists} rule with the existential only in the precondition (\texttt{triple\_hexists}), and
derives the form used here from it (\texttt{triple\_hexists2}).

% \end{document}
